\nonstopmode{}
\documentclass[a4paper]{book}
\usepackage[times,inconsolata,hyper]{Rd}
\usepackage{makeidx}
\makeatletter\@ifl@t@r\fmtversion{2018/04/01}{}{\usepackage[utf8]{inputenc}}\makeatother
% \usepackage{graphicx} % @USE GRAPHICX@
\makeindex{}
\begin{document}
\chapter*{}
\begin{center}
{\textbf{\huge Package `rcamelspe'}}
\par\bigskip{\large \today}
\end{center}
\ifthenelse{\boolean{Rd@use@hyper}}{\hypersetup{pdftitle = {rcamelspe: Catchment Attributes and Meteorology for Large-Sample Studies - Peru (CAMELS-PE) Data Management}}}{}
\ifthenelse{\boolean{Rd@use@hyper}}{\hypersetup{pdfauthor = {Harold Llauca}}}{}
\begin{description}
\raggedright{}
\item[Title]\AsIs{Catchment Attributes and Meteorology for Large-Sample Studies -
Peru (CAMELS-PE) Data Management}
\item[Version]\AsIs{0.1.0}
\item[Description]\AsIs{Tools to efficiently download, load, manage, and filter the CAMELS-PE dataset for hydroclimatic and hydrological analysis in Peru. Uses 'arrow' and 'collapse' for high-performance reading and manipulation of large datasets.}
\item[License]\AsIs{MIT + file LICENSE}
\item[Encoding]\AsIs{UTF-8}
\item[Roxygen]\AsIs{list(markdown = TRUE)}
\item[Imports]\AsIs{arrow, collapse (>= 2.0.0), sf, tidyselect}
\item[Suggests]\AsIs{testthat (>= 3.0.0)}
\item[Config/testthat/edition]\AsIs{3}
\item[Config/roxygen2/version]\AsIs{8.0.0}
\item[RoxygenNote]\AsIs{7.3.3}
\item[NeedsCompilation]\AsIs{no}
\item[Author]\AsIs{Harold Llauca [aut, cre]}
\item[Maintainer]\AsIs{Harold Llauca }\email{hllauca@senamhi.gob.pe}\AsIs{}
\end{description}
\Rdcontents{Contents}
\HeaderA{download\_pe\_data}{Download the CAMELS-PE dataset from Zenodo}{download.Rul.pe.Rul.data}
%
\begin{Description}
Downloads the CAMELS-PE dataset zip file from the Zenodo repository
and optionally unzips it in the destination folder.
\end{Description}
%
\begin{Usage}
\begin{verbatim}
download_pe_data(
  dest_dir = "data-raw",
  unzip = TRUE,
  overwrite = FALSE,
  quiet = FALSE
)
\end{verbatim}
\end{Usage}
%
\begin{Arguments}
\begin{ldescription}
\item[\code{dest\_dir}] Character. Destination directory where the dataset should be saved.
Defaults to "data-raw".

\item[\code{unzip}] Logical. Should the downloaded file be unzipped? Default is \code{TRUE}.

\item[\code{overwrite}] Logical. If the destination files already exist, should they be overwritten?
Default is \code{FALSE}.

\item[\code{quiet}] Logical. Should download progress be suppressed? Default is \code{FALSE}.
\end{ldescription}
\end{Arguments}
%
\begin{Value}
Character. The directory path where the dataset is located or extracted.
\end{Value}
%
\begin{Examples}
\begin{ExampleCode}
## Not run: 
# Download and unzip to a local folder
download_pe_data(dest_dir = "data-raw")

## End(Not run)
\end{ExampleCode}
\end{Examples}
\HeaderA{get\_camels\_pe\_path}{Get the path to the CAMELS-PE dataset}{get.Rul.camels.Rul.pe.Rul.path}
%
\begin{Description}
Retrieves the path to the CAMELS-PE dataset. It looks up:
\begin{enumerate}

\item{} The package option \code{rcamelspe.path}.
\item{} The environment variable \code{CAMELS\_PE\_PATH}.
\item{} Default search paths in the current working directory (\code{raw-data/CAMELS-PE}, \code{data-raw/camels-pe}, etc.).

\end{enumerate}

\end{Description}
%
\begin{Usage}
\begin{verbatim}
get_camels_pe_path()
\end{verbatim}
\end{Usage}
%
\begin{Value}
Character path or \code{NULL} if not found.
\end{Value}
%
\begin{Examples}
\begin{ExampleCode}
get_camels_pe_path()
\end{ExampleCode}
\end{Examples}
\HeaderA{load\_pe\_attributes}{Load CAMELS-PE catchment attributes}{load.Rul.pe.Rul.attributes}
%
\begin{Description}
Loads one or more catchment attribute files and merges them by \code{gauge\_id}.
\end{Description}
%
\begin{Usage}
\begin{verbatim}
load_pe_attributes(attributes = "all", path = get_camels_pe_path())
\end{verbatim}
\end{Usage}
%
\begin{Arguments}
\begin{ldescription}
\item[\code{attributes}] Character vector. The attributes to load. Can be any combination of
\code{"topographic"}, \code{"climatic"}, \code{"geologic"}, \code{"soil"}, \code{"landcover"}, \code{"intervention"},
and \code{"signatures"}, or \code{"all"} (default) to load and merge all attributes.

\item[\code{path}] Character. Path to the CAMELS-PE dataset directory. If NULL,
retrieved via \code{\LinkA{get\_camels\_pe\_path()}{get.Rul.camels.Rul.pe.Rul.path}}.
\end{ldescription}
\end{Arguments}
%
\begin{Value}
A data frame containing the merged attributes.
\end{Value}
%
\begin{Examples}
\begin{ExampleCode}
## Not run: 
# Load all attributes merged
attrs_all <- load_pe_attributes()

# Load only topographic and climatic attributes
attrs_sub <- load_pe_attributes(c("topographic", "climatic"))

## End(Not run)
\end{ExampleCode}
\end{Examples}
\HeaderA{load\_pe\_geospatial}{Load CAMELS-PE geospatial data}{load.Rul.pe.Rul.geospatial}
%
\begin{Description}
Loads the geospatial catchment boundaries (polygons) or gauging station locations (points)
from the CAMELS-PE dataset.
\end{Description}
%
\begin{Usage}
\begin{verbatim}
load_pe_geospatial(
  type = c("catchments", "gauges"),
  gauge_ids = NULL,
  path = get_camels_pe_path()
)
\end{verbatim}
\end{Usage}
%
\begin{Arguments}
\begin{ldescription}
\item[\code{type}] Character. Either \code{"catchments"} to load catchment boundary polygons
or \code{"gauges"} to load gauge point locations.

\item[\code{gauge\_ids}] Character vector. Station identifiers to filter the spatial features.
If \code{NULL} (default), loads all features.

\item[\code{path}] Character. Path to the CAMELS-PE dataset directory. If NULL,
retrieved via \code{\LinkA{get\_camels\_pe\_path()}{get.Rul.camels.Rul.pe.Rul.path}}.
\end{ldescription}
\end{Arguments}
%
\begin{Value}
An \code{sf} spatial object containing the requested geometries.
\end{Value}
%
\begin{Examples}
\begin{ExampleCode}
## Not run: 
# Load all catchment boundaries
catchments <- load_pe_geospatial(type = "catchments")

# Load specific gauging stations
gauges_sub <- load_pe_geospatial(type = "gauges", gauge_ids = c("PE_110139", "PE_111151"))

## End(Not run)
\end{ExampleCode}
\end{Examples}
\HeaderA{load\_pe\_metadata}{Load CAMELS-PE metadata}{load.Rul.pe.Rul.metadata}
%
\begin{Description}
Loads the gauging station metadata or the data dictionary from the CAMELS-PE dataset.
\end{Description}
%
\begin{Usage}
\begin{verbatim}
load_pe_metadata(
  type = c("stations", "dictionary"),
  path = get_camels_pe_path()
)
\end{verbatim}
\end{Usage}
%
\begin{Arguments}
\begin{ldescription}
\item[\code{type}] Character. Either \code{"stations"} to load the gauging station metadata
or \code{"dictionary"} to load the data dictionary.

\item[\code{path}] Character. Path to the CAMELS-PE dataset directory. If NULL,
retrieved via \code{\LinkA{get\_camels\_pe\_path()}{get.Rul.camels.Rul.pe.Rul.path}}.
\end{ldescription}
\end{Arguments}
%
\begin{Value}
A data frame (tibble-like) containing the requested metadata.
\end{Value}
%
\begin{Examples}
\begin{ExampleCode}
## Not run: 
# Load stations metadata
stations <- load_pe_metadata(type = "stations")

# Load data dictionary
data_dict <- load_pe_metadata(type = "dictionary")

## End(Not run)
\end{ExampleCode}
\end{Examples}
\HeaderA{load\_pe\_timeseries}{Load CAMELS-PE daily timeseries data}{load.Rul.pe.Rul.timeseries}
%
\begin{Description}
Efficiently loads the daily hydroclimatic timeseries for Peruvian catchments.
Optimized to load either specific catchments from individual files or the entire
dataset using 'arrow' and 'collapse'.
\end{Description}
%
\begin{Usage}
\begin{verbatim}
load_pe_timeseries(
  gauge_ids = NULL,
  variables = NULL,
  path = get_camels_pe_path(),
  parse_dates = TRUE,
  use_arrow = TRUE
)
\end{verbatim}
\end{Usage}
%
\begin{Arguments}
\begin{ldescription}
\item[\code{gauge\_ids}] Character vector. Station/catchment identifiers (e.g. \code{"PE\_110139"}).
If \code{NULL} (default), loads timeseries for all 136 catchments.

\item[\code{variables}] Character vector. Names of the hydroclimatic variables to select.
If \code{NULL} (default), all variables are loaded. Allowed variables:
\code{"prec"}, \code{"prec\_var"}, \code{"flow\_obs"}, \code{"flow\_sim"}, \code{"pet"}, \code{"tmin"},
\code{"tmean"}, \code{"tmax"}, \code{"srad"}, \code{"vprp"}.

\item[\code{path}] Character. Path to the CAMELS-PE dataset directory. If NULL,
retrieved via \code{\LinkA{get\_camels\_pe\_path()}{get.Rul.camels.Rul.pe.Rul.path}}.

\item[\code{parse\_dates}] Logical. Should the \code{date} column be parsed into R Date objects?
Default is \code{TRUE}.

\item[\code{use\_arrow}] Logical. Should the 'arrow' package be used for reading?
Default is \code{TRUE} (recommended for speed).
\end{ldescription}
\end{Arguments}
%
\begin{Value}
A data frame containing the daily timeseries.
\end{Value}
%
\begin{Examples}
\begin{ExampleCode}
## Not run: 
# Load timeseries for a specific station
ts_station <- load_pe_timeseries("PE_110139")

# Load only precipitation and observed streamflow for all stations
ts_sub <- load_pe_timeseries(variables = c("prec", "flow_obs"))

## End(Not run)
\end{ExampleCode}
\end{Examples}
\HeaderA{set\_camels\_pe\_path}{Set the path to the CAMELS-PE dataset}{set.Rul.camels.Rul.pe.Rul.path}
%
\begin{Description}
Sets the path to the directory containing the CAMELS-PE dataset.
The path is stored in the package options.
\end{Description}
%
\begin{Usage}
\begin{verbatim}
set_camels_pe_path(path)
\end{verbatim}
\end{Usage}
%
\begin{Arguments}
\begin{ldescription}
\item[\code{path}] Character. Path to the CAMELS-PE dataset directory.
\end{ldescription}
\end{Arguments}
%
\begin{Value}
Invisible NULL.
\end{Value}
%
\begin{Examples}
\begin{ExampleCode}
set_camels_pe_path("data-raw/CAMELS-PE")
\end{ExampleCode}
\end{Examples}
\printindex{}
\end{document}
